% Upstream source: ~/mm/analysis/hp_error_rate/hp_error_rate_report.tex
% This copy is distributed alongside speconsense for convenient access; edits
% should be made upstream and synced here, not the other way around.
\documentclass[11pt]{article}

\usepackage[margin=1in]{geometry}
\usepackage{amsmath,amssymb}
\usepackage{graphicx}
\usepackage{booktabs}
\usepackage{hyperref}
\usepackage[numbers]{natbib}
\usepackage{textcomp}
\usepackage{multirow}
\usepackage{lineno}

\title{Homopolymer error rate analysis for nanopore amplicon sequencing}
\author{Josh Walker}
\date{Internal review draft, May 2026 (commit 5d33d1d)}

\begin{document}
\maketitle
\linenumbers

\begin{abstract}
We quantify homopolymer (hp) length error rates in Oxford Nanopore amplicon
sequencing by base and run length, using two independent sequencing runs of
fungal ITS amplicons spanning two basecaller generations (Dorado SUP v3.5 and
v5.0) on R10.4.1 chemistry. After filtering biological variants via
statistical outlier and bimodal distribution detection, we find that implied
per-position error rates at hp lengths 1--5 are consistently below 1.4\% for
Dorado v5.0 and below 2.3\% for v3.5---comparable to the non-hp substitution
baseline. Error rates rise steeply at lengths 6+, with G/C 2--4$\times$ worse
than A/T. We also investigate substitution errors at hp boundaries (interface
positions) and find no evidence of excess artifactual variation: well-supported
variants at interface positions occur at or below the rate of non-hp variants
(rate ratios of 1.03 [95\% CI: 0.88--1.19] for v3.5 and 0.88 [0.79--0.99] for
v5.0). These results support extending
speconsense's critical error rate (CER) framework to evaluate hp-position
variants at lengths 1--5 statistically, rather than discarding them
unconditionally via blanket homopolymer normalization. The CER extension should
operate at the run level, using empirical error distributions by (base, length)
context; the per-position rate serves as a compact parameterization for
cross-length comparison, not a mechanistic model of per-base independence.
\end{abstract}

\section{Introduction}

\subsection{Homopolymer errors in nanopore sequencing}

Homopolymer (hp) length errors are the dominant systematic error mode in
nanopore sequencing. Approximately 47\% of all nanopore sequencing errors
originate in homopolymer regions~\cite{delahaye2021}, and even with the latest
chemistry and basecalling, homopolymer-length indels account for 56\% of
remaining errors in state-of-the-art genome assemblies~\cite{wick2025}. The
underlying mechanism is well understood: nanopore sequencing infers base
identity from ionic current changes as DNA translocates through the pore, but
consecutive identical bases produce no current change, making the length of a
homopolymer run difficult to determine from signal alone.

While the qualitative understanding that nanopore sequencing struggles with
homopolymers is widespread, quantitative characterization of error rates by
homopolymer length and base has been limited. The most comprehensive published
data comes from two studies:

\begin{itemize}
\item \textbf{Sereika et al.\ (2022)}~\cite{sereika2022} characterized
  per-length, per-base hp accuracy at both read and assembly level using a Zymo
  mock community, with Illumina-polished PacBio HiFi assemblies as ground
  truth. Their comparison of R9.4.1 and R10.4 chemistry demonstrated a
  20--30~percentage point improvement in read-level accuracy at hp lengths 7--9
  with the newer dual reader head.

\item \textbf{Ni et al.\ (2023)}~\cite{ni2023} independently confirmed these
  trends on human whole genome data, benchmarking R9.4.1 and R10.4 against a
  T2T reference genome.
\end{itemize}

Both studies focused primarily on hp lengths $\geq 5$, where errors are
dramatic and consequential for genome assembly. Short homopolymers (lengths
1--4) were noted to be highly accurate ({>}95\% at read level) but were not
characterized in detail, as they were not considered problematic for typical
applications.

The rapid pace of improvement in nanopore basecalling further limits the
applicability of published data. Sereika's R10.4 results used Guppy SUP
basecalling; our data spans two basecaller generations (Dorado SUP v3.5.x and
v5.0.0) on R10.4.1 chemistry. No published study provides per-length hp error
rates for either configuration.

\subsection{The speconsense problem}

This gap in quantitative understanding has practical consequences for nanopore
amplicon analysis. Speconsense's variant phasing pipeline uses homopolymer
normalization to avoid calling false variants caused by hp length errors.
Currently, \emph{all} homopolymer length differences are treated as
insignificant---a difference between \texttt{ATC} and \texttt{ATTC} is treated
identically to a difference between \texttt{AAAAAAAAA} and
\texttt{AAAAAAAAAA}. This is conservative but lossy: it discards variants that
differ only at homopolymer positions, even when those differences are very
likely to be biological.

The loss is not merely theoretical. At short hp lengths (1--4), the published
literature suggests read-level accuracy exceeds 95\%, meaning that a length
difference at these positions is far more likely to reflect real biology than
sequencing error. Yet speconsense discards this information.

Speconsense's critical error rate (CER) framework~\cite{walker2026} provides a
statistical approach for evaluating whether an observed variant is plausibly
explained by sequencing error, but it currently applies only to non-hp
positions. Extending the CER framework to hp positions requires quantitative
error rate estimates by hp length and base---exactly the data that is missing
from the published literature for current chemistry. The companion
report~\cite{cer_practice} integrates the rates derived in this analysis into
a context-aware CER framework and its implementation in speconsense; the
present report provides the empirical foundation.

\subsection{This analysis}

This analysis quantifies nanopore hp error rates by homopolymer length and
base, using real amplicon data from two independent sequencing runs spanning
two basecaller generations (Dorado SUP v3.5 and v5.0), to determine whether
the CER framework can be extended to evaluate hp-position variants
statistically rather than discarding them unconditionally. We also investigate
whether substitution errors at homopolymer boundaries (interface positions)
require special treatment, or can be evaluated using the same framework as
non-hp positions.


\section{Methods}

\subsection{Data}

We analyzed specimens from sequencing runs of fungal ITS amplicons
(${\sim}700$~bp), all produced by Mycota Lab (Hoosier Mushroom Society)
following the protocol of Russell~\cite{russell2023} and sequenced on
Oxford Nanopore R10.4.1 E8.2 chemistry. Reads were demultiplexed with
specimux and clustered with speconsense using
\texttt{-{}-presample~0 -{}-max-sample-size~10000} to avoid quality-sorted
presampling bias.

\begin{itemize}
\item \textbf{ont37}: 94 qualifying specimens from 651 with c1 MSAs
  (environmental specimens). Basecalled with Dorado SUP v3.5.1 (70\% of reads)
  and v3.5.2 260bps (30\%), sequenced February 2024. Two sample IDs (ONT037
  and ONT37B) were combined; both contain reads from both basecaller models.
\item \textbf{ont98}: 245 qualifying specimens from 837 total (environmental
  specimens). Basecalled with Dorado SUP v5.0.0, sequenced August 2025.
\end{itemize}

For ont98, qualifying specimens had a primary cluster with $\geq 200$ reads
representing $\geq 70\%$ of total reads. For ont37, thresholds were relaxed
to $\geq 100$ reads and $\geq 50\%$ dominance due to lower overall sequencing
depth (median c1 read count of 59 vs 200+ for ont98).

\subsection{HP length extraction}

For each specimen, the SPOA multiple sequence alignment (MSA) was parsed to
extract called hp lengths at each hp run in the consensus. The extraction
algorithm:

\begin{enumerate}
\item Identify maximal hp runs in the ungapped consensus sequence.
\item For each run, locate the flanking non-hp bases (``anchors'') in the MSA.
\item For each read, count consecutive instances of the hp base between the
  anchors in the gapped alignment. This gives the read's called hp length for
  that run.
\end{enumerate}

\subsection{Ground truth determination}

Rather than using the SPOA consensus hp length as ground truth (which can be
biased by SPOA's heaviest-bundle algorithm), we used the \textbf{empirical
mode} of called hp lengths across all reads at each position. At depths of
200+ reads on R10.4.1 SUP, the mode is the correct length for hp $\leq
{\sim}9$~\cite{sereika2022}.

\subsection{Error classification}

At each hp run, reads were classified as:

\begin{itemize}
\item \textbf{Correct}: called length equals the mode
\item \textbf{HP deletion ($-k$)}: called length is $k$ bases shorter than the mode
\item \textbf{HP insertion ($+k$)}: called length is $k$ bases longer than the mode
\item \textbf{Interface substitution}: a base within the hp region is replaced
  by an adjacent flanking base (see Section~\ref{sec:interface})
\item \textbf{Other substitution}: a non-flanking base appears within the hp region
\item \textbf{Complex}: multiple error types within the hp region
\end{itemize}

Non-hp positions (singletons not adjacent to any run) were analyzed separately
as a calibration baseline.

\subsection{Scope of measurement}

The error rates reported here are not raw nanopore sequencing error rates.
They are \textbf{effective error rates at the post-clustering stage} of the
speconsense pipeline: reads have been demultiplexed, clustered by MCL, and
aligned with SPOA before analysis. Speconsense's \texttt{-{}-min-identity}
parameter controls a quality filter during clustering that excludes reads
below a pairwise identity threshold to the cluster consensus. Reads that fail
this filter are never included in the MSA and therefore do not contribute to
the error rates we measure.

This is the correct quantity for parameterizing the CER framework, which
operates on the same post-clustering MSA. The effective error rate is what
determines whether an observed variant in a speconsense cluster is plausibly
explained by sequencing error. However, it means that our error rates are
conditioned on speconsense's clustering parameters and are not directly
comparable to raw read-level error rates reported in the literature.

To assess the sensitivity of our results to this filtering, we repeated the
analysis on ont37 and ont98 with \texttt{-{}-min-identity~0.85} (vs the
default of 0.90), which admits more divergent reads into clusters. Results are
reported in Section~\ref{sec:identity_sensitivity}.


\section{Results (ont98)}

\subsection{Dataset summary}

\begin{table}[h]
\centering
\begin{tabular}{lr}
\toprule
Metric & Value \\
\midrule
Specimens analyzed & 245 \\
Total HP runs & 129,673 \\
Total read-level HP observations & 83,429,870 \\
Non-HP position observations & 60,059,474 \\
\bottomrule
\end{tabular}
\end{table}

\subsection{Non-HP error rates (calibration)}

\begin{table}[h]
\centering
\begin{tabular}{lr}
\toprule
Error Type & Rate \\
\midrule
Substitution & 0.59\% \\
Deletion & 0.40\% \\
Insertion & 0.68\% \\
\textbf{Total (sub + del)} & \textbf{0.99\%} \\
\bottomrule
\end{tabular}
\end{table}

These rates are consistent with R10.4.1 SUP read accuracy of ${\sim}99\%$,
confirming the calibration baseline.

\subsection{HP error rates by base and length}

Table~\ref{tab:hp_error} shows the fraction of reads calling the correct hp
length (mode), and the fraction with $\pm 1$ or larger deviations. Only
contexts with $\geq 100$ observations are shown.

\begin{table}[h]
\centering
\small
\caption{HP error rates by base and length (ont98, raw).}
\label{tab:hp_error}
\begin{tabular}{llrrrrrr}
\toprule
Base & Length & N runs & N obs & Correct & Del ($-1$) & Ins ($+1$) & $\pm 2{+}$ \\
\midrule
\multirow{5}{*}{A}
  & 1 & 20,114 & 12,930,648 & 99.32\% & 0.55\% & 0.12\% & 0.01\% \\
  & 2 & 8,479 & 5,487,156 & 98.25\% & 1.45\% & 0.21\% & 0.10\% \\
  & 3 & 1,644 & 1,067,803 & 97.26\% & 2.13\% & 0.29\% & 0.32\% \\
  & 4 & 595 & 378,225 & 96.15\% & 2.44\% & 0.41\% & 1.00\% \\
  & 5 & 93 & 56,312 & 92.92\% & 4.57\% & 1.03\% & 1.48\% \\
\midrule
\multirow{5}{*}{T}
  & 1 & 24,719 & 15,845,528 & 99.30\% & 0.54\% & 0.15\% & 0.01\% \\
  & 2 & 7,337 & 4,770,094 & 98.05\% & 1.60\% & 0.24\% & 0.10\% \\
  & 3 & 2,927 & 1,852,010 & 96.74\% & 2.53\% & 0.37\% & 0.37\% \\
  & 4 & 657 & 430,793 & 95.60\% & 3.03\% & 0.60\% & 0.77\% \\
  & 5 & 248 & 163,547 & 93.70\% & 3.89\% & 0.99\% & 1.42\% \\
\midrule
\multirow{5}{*}{C}
  & 1 & 24,848 & 15,986,540 & 99.21\% & 0.69\% & 0.09\% & 0.00\% \\
  & 2 & 4,770 & 3,053,272 & 98.44\% & 1.16\% & 0.26\% & 0.14\% \\
  & 3 & 999 & 647,920 & 96.92\% & 2.15\% & 0.53\% & 0.40\% \\
  & 4 & 309 & 189,751 & 95.47\% & 3.00\% & 0.85\% & 0.68\% \\
  & 5 & 55 & 36,493 & 90.54\% & 5.60\% & 1.87\% & 1.99\% \\
\midrule
\multirow{5}{*}{G}
  & 1 & 23,491 & 15,095,593 & 99.31\% & 0.60\% & 0.08\% & 0.01\% \\
  & 2 & 6,858 & 4,476,352 & 98.38\% & 1.25\% & 0.21\% & 0.16\% \\
  & 3 & 923 & 576,762 & 96.74\% & 2.26\% & 0.50\% & 0.50\% \\
  & 4 & 245 & 147,437 & 93.99\% & 4.25\% & 0.81\% & 0.96\% \\
  & 5 & 132 & 86,883 & 91.01\% & 5.58\% & 1.73\% & 1.68\% \\
\bottomrule
\end{tabular}
\end{table}

Extended results for hp lengths 6--14 (all bases) are provided in the
supplementary tables. Key trends at longer lengths: accuracy drops to
67--91\% at length~7, 63--86\% at length~8, and below 62\% at length~10.

\subsection{Implied per-position error rate}

The run-level fraction correct compounds over the length of the hp run. To
compare error rates across different hp lengths and against the non-hp
baseline, we compute the implied per-position error rate $p$ such that
$(1-p)^L$ equals the observed fraction correct. This decomposition is a
convenience for cross-length comparison, not a mechanistic model: hp length
errors are signal-level events affecting the entire run at once, not
independent per-position errors. The per-position rate should be understood as
a compact parameterization of the run-level error probability, useful for
normalizing out the trivial effect of run length on aggregate accuracy.

\begin{table}[h]
\centering
\caption{Pooled per-position error rates by HP length (ont98, raw).}
\label{tab:per_pos}
\begin{tabular}{rrrr}
\toprule
Length & Run Error & Per-pos Rate & vs Non-HP (0.99\%) \\
\midrule
1 & 0.72\% & 0.72\% & $0.72\times$ \\
2 & 1.74\% & 0.87\% & $0.88\times$ \\
3 & 3.10\% & 1.04\% & $1.05\times$ \\
4 & 4.45\% & 1.13\% & $1.14\times$ \\
5 & 7.45\% & 1.54\% & $1.55\times$ \\
6 & 11.03\% & 1.93\% & $1.95\times$ \\
7 & 19.22\% & 3.00\% & $3.03\times$ \\
8 & 28.92\% & 4.18\% & $4.22\times$ \\
9 & 33.01\% & 4.35\% & $4.40\times$ \\
10 & 38.13\% & 4.69\% & $4.73\times$ \\
\bottomrule
\end{tabular}
\end{table}

\textbf{Key finding}: At hp lengths 1--4, the implied per-position error rate
is comparable to or below the non-hp baseline (0.7--1.1\% vs 0.99\%). This
means that the run-level error probability at these lengths is within the range
that would be expected if the per-position error rate were similar to the
non-hp baseline---i.e., short hp runs are not inherently noisier than non-hp
positions, once the effect of run length is accounted for. At lengths 5+, the
implied per-position rate rises above the baseline, reaching
2--5$\times$ at lengths 7--9, reflecting the increasing difficulty of
signal-level hp length estimation. (This per-position decomposition is a
comparison metric, not a claim about the error mechanism; the full run-level
error distributions are in Table~\ref{tab:hp_error} and in the supplementary
\texttt{hp\_distributions.tsv} tables.)

\subsection{Qualitative observations}

\begin{enumerate}
\item \textbf{Short hp runs (length 1--4) have very low error rates.} At
  length~1 (singletons), the hp ``error'' rate of 0.55--0.69\% is comparable to
  the non-hp baseline of 0.99\%. At length~4, accuracy is still 94--96\%. These
  positions are strong candidates for CER-based variant recovery.

\item \textbf{Accuracy degrades monotonically with length, with a steep drop at
  length 7--8.} At length~7, accuracy ranges from 67\% (G) to 91\% (T). By
  length~10, accuracy is 54--62\%---below the threshold where the mode is
  reliably the correct length.

\item \textbf{G/C homopolymers are harder than A/T.} G is consistently worst
  across all lengths. At length~5: A~92.9\%, T~93.7\%, C~90.5\%, G~91.0\%.
  The gap widens at longer lengths.

\item \textbf{Deletions dominate over insertions at all lengths.} The
  deletion\,:\,insertion ratio is 3--5\,:\,1 for moderate hp lengths (4--7),
  consistent with the known nanopore bias toward underestimating hp length.

\item \textbf{Our error rates are substantially better than Sereika et
  al.~(2022).} At length~8, Sereika reported ${\sim}$45--55\% correct for
  R10.4 A/T reads; we observe 68--86\%. This improvement likely reflects the
  newer basecaller (Dorado SUP v5.0.0 vs Guppy SUP) and possibly chemistry
  refinements (R10.4.1 vs R10.4), though the comparison is not exact: our
  effective error rates are conditioned on speconsense's post-clustering MSA
  (Section~2.5), whereas Sereika measured raw read-level accuracy against a
  known reference.
\end{enumerate}

\subsection{SPOA consensus bias}

Only 9 out of 129,673 hp runs (0.007\%) showed a discrepancy between the SPOA
consensus hp length and the empirical mode. The heaviest-bundle bias is
negligible in this dataset, likely because the read depths are high enough that
the mode strongly dominates.

\subsection{Candidate biological hp variants}

The bimodal distribution detector flagged 106 positions (0.08\% of all hp runs)
where a second hp length was present at rates significantly above the error
model prediction. Many of these show approximately 50:50 splits consistent with
heterozygous positions---variants that speconsense currently discards.


\section{Interface Homopolymers}
\label{sec:interface}

\subsection{Definition}

An ``interface homopolymer'' substitution occurs at the boundary between a
homopolymer run and an adjacent base, where the substitution produces a
sequence that can be reinterpreted as a hp length change on both adjacent runs.
For example:

\begin{verbatim}
    Consensus: ...TT|CC|TT...
    Variant:   ...TT|CT|TT...  (C->T at position 2 of the CC run)
\end{verbatim}

This can equivalently be read as: the T-run extended from 2 to 3, and the
C-run shortened from 2 to 1. The adjusted-identity library treats this as
hp-equivalent (0~edits), so speconsense does not detect it as a variant.

\subsection{Prevalence}

We classified every read observation at hp runs (length $\geq 2$) by error
type:

\begin{table}[h]
\centering
\begin{tabular}{lrr}
\toprule
Error Type & Read Observations & Rate \\
\midrule
Correct & 22,585,613 & 95.77\% \\
Pure HP deletion & 249,385 & 1.06\% \\
Pure HP insertion & ${\sim}$100,000 & 0.42\% \\
\textbf{Interface substitution} & \textbf{228,035} & \textbf{0.97\%} \\
Other substitution & 122,236 & 0.52\% \\
Complex & 47,000 & 0.20\% \\
\bottomrule
\end{tabular}
\end{table}

Interface hp substitutions occur at \textbf{nearly the same rate as pure hp
length errors} (0.97\% vs 1.06\%).

\subsection{Expanded analysis: interface variant rate vs non-HP baseline}

To determine whether interface hp substitutions are primarily signal-level
artifacts or biological variants, we conducted an expanded analysis across 663
specimens (all specimens with $\geq 100$ total reads, removing the original
selection bias toward low-variation specimens).

\textbf{Methodology}: For each specimen, reads from all same-organism clusters
($\geq 90\%$ pairwise identity to the primary cluster consensus) were merged
and re-aligned with SPOA to produce a combined MSA. This ensures that
biological non-HP substitutions---which speconsense phases into separate
clusters---are visible alongside interface HP substitutions, enabling a fair
comparison. Each consensus position was classified as non-HP, HP interface
(first or last base of a run of length $\geq 2$ adjacent to a different base),
or HP interior.

The key metric is the \textbf{fraction of positions showing a well-supported
variant} (minor allele frequency $\geq$ threshold), compared across position
classes. If interface positions produce artifactual variants due to
signal-level boundary confusion, they should show an excess of well-supported
variants relative to non-HP positions.

\begin{table}[h]
\centering
\caption{Variant position rates by position class (ont98, 663 specimens).}
\label{tab:interface_ont98}
\begin{tabular}{lrrrl}
\toprule
MAF Threshold & Non-HP & HP Interface & HP Interior & Ratio \\
\midrule
$\geq 10\%$ & 0.319\% & 0.276\% & 0.328\% & 0.87 \\
$\geq 20\%$ & 0.265\% & 0.231\% & 0.259\% & 0.87 \\
$\geq 30\%$ & 0.219\% & 0.184\% & 0.199\% & 0.84 \\
$\geq 40\%$ & 0.158\% & 0.133\% & 0.132\% & 0.84 \\
\bottomrule
\end{tabular}
\end{table}

The rate ratio of interface to non-HP variant positions at MAF $\geq 20\%$ is
\textbf{0.880 [95\% bootstrap CI: 0.785--0.988]}, with the confidence interval
entirely below 1.0. Interface positions produce well-supported variants at a
slightly \emph{lower} rate than non-HP positions.

Of the 457 interface positions with MAF $\geq 20\%$, only 173 (38\%) had an
interface-compatible top allele (i.e., the variant base matched the adjacent
run's base). This is close to the 33\% expected by chance when there are three
possible alternative bases, indicating no enrichment of boundary-confusion
substitutions among well-supported variants.

\textbf{Conclusion}: There is no evidence that HP interface positions generate
artifactual well-supported variants at rates above the non-HP baseline.


\section{Biological variant filtering}
\label{sec:bio_filtering}

The raw error rates above include HP positions where the observed variation is
biological (genuine length polymorphisms) rather than sequencing error. The
filtered analyses in subsequent sections (and the re-estimation in
Section~\ref{sec:reestimation}) identify and exclude these positions using two
detectors:

\begin{itemize}
\item \textbf{Statistical outliers}: HP positions where the error rate
  significantly exceeds the pooled average for that context (one-sided binomial
  test, Bonferroni-corrected $\alpha = 0.01$).
\item \textbf{Bimodal positions}: HP positions where a second allele is present
  at ${>}\,2\times$ the expected error rate and ${>}\,10\%$ of reads.
\end{itemize}

For ont98, the detectors flag 1,656 outlier positions and 106 bimodal positions
(1,657 total excluded, 1.3\% of HP runs). The filtered per-position rates are
referenced as ``ont98 filtered'' in the cross-basecaller comparison
(Section~\ref{sec:older_basecaller}) and as the input distribution for
re-estimation (Section~\ref{sec:reestimation}).


\section{Older Basecaller Comparison (ont37)}
\label{sec:older_basecaller}

\subsection{Dataset}

To assess the effect of basecaller version on HP error rates, we analyzed
ont37, a sequencing run from February~2024 basecalled with Dorado SUP v3.5.x
(vs v5.0.0 for ont98). The data was combined from two sample IDs (ONT037 and
ONT37B), with 70\% of reads basecalled using the standard SUP v3.5.1 model
and 30\% using the 260bps SUP v3.5.2 model. Both datasets use R10.4.1 E8.2
chemistry.

Due to lower sequencing depth (median c1 read count of 59), qualifying
thresholds were relaxed to $\geq 100$ primary reads and $\geq 50\%$ dominance,
yielding 94 qualifying specimens from 651 with c1 MSAs. Total filtered HP
observations: 10.2~million (vs 82.0~million for ont98).

\subsection{Non-HP error rates}

\begin{table}[h]
\centering
\caption{Non-HP error rates across both datasets.}
\label{tab:nonhp_all}
\begin{tabular}{lrrrr}
\toprule
Dataset & Substitution & Deletion & Insertion & Total \\
\midrule
ont37 (v3.5) & 1.44\% & 0.99\% & 1.54\% & 3.97\% \\
ont98 (v5.0) & 0.59\% & 0.40\% & 0.68\% & 1.67\% \\
\bottomrule
\end{tabular}
\end{table}

The ont37 non-HP total error rate of 3.97\% is 2.4$\times$ higher than ont98.
All error classes are elevated: substitutions 2.4$\times$, deletions
2.5$\times$, and insertions 2.3$\times$.

\subsection{HP error rates and biological variant filtering}

After biological variant filtering (817 outlier positions, 664 bimodal
positions, 1,092 total excluded = 2.2\%), the filtered per-position HP error
rates across both datasets are:

\begin{table}[h]
\centering
\caption{Filtered per-position HP error rates: cross-basecaller comparison.}
\label{tab:three_run}
\begin{tabular}{rrrr}
\toprule
HP Len & ont37 (v3.5) & ont98 (v5.0) & ont37/ont98 \\
\midrule
1 & 1.56\% & 0.67\% & 2.32 \\
2 & 1.84\% & 0.83\% & 2.21 \\
3 & 2.17\% & 0.97\% & 2.23 \\
4 & 2.27\% & 0.99\% & 2.30 \\
5 & 2.13\% & 1.13\% & 1.89 \\
6 & 2.90\% & 1.57\% & 1.85 \\
7 & 4.21\% & 2.23\% & 1.89 \\
\bottomrule
\end{tabular}
\end{table}

At L=1--5, ont37 error rates are consistently ${\sim}$2.2$\times$ those of
ont98, demonstrating substantial improvement from Dorado v3.5 to v5.0.
Run-to-run variation within a basecaller version (independent of this
v3.5$\to$v5.0 effect) is expected but not characterized in this report.

\subsection{Whole-run accuracy by HP length}

\begin{table}[h]
\centering
\caption{Whole-run accuracy (fraction correct) across datasets (filtered).}
\label{tab:frac_correct_all}
\begin{tabular}{rrr}
\toprule
HP Length & ont37 (v3.5) & ont98 (v5.0) \\
\midrule
1 & 98.4\% & 99.3\% \\
2 & 96.4\% & 98.3\% \\
3 & 93.6\% & 97.1\% \\
4 & 91.2\% & 96.1\% \\
5 & 89.8\% & 94.5\% \\
6 & 83.8\% & 91.0\% \\
7 & 74.0\% & 85.4\% \\
\bottomrule
\end{tabular}
\end{table}

Even with the older v3.5 basecaller, short HP accuracy remains high: 98.4\% at
L=1 and 91.2\% at L=4. The CER framework can be applied to short HPs across
basecaller versions, with the per-position error rate parameterized
appropriately for each.

\subsection{Base composition (ont37)}

\begin{table}[h]
\centering
\caption{Per-base variation in filtered HP error rates (ont37).}
\label{tab:base_ont37}
\begin{tabular}{rlllr}
\toprule
HP Length & A & C & G & T \\
\midrule
1 & 1.26\% & 1.95\% & 1.71\% & 1.31\% \\
2 & 1.79\% & 1.63\% & 1.82\% & 2.02\% \\
3 & 1.61\% & 3.00\% & 2.35\% & 2.23\% \\
4 & 2.32\% & 3.11\% & 3.17\% & 1.92\% \\
5 & 2.11\% & 3.76\% & 2.37\% & 1.88\% \\
\bottomrule
\end{tabular}
\end{table}

The base composition pattern in ont37 shows larger C/G vs A/T divergence at
shorter lengths compared to the v5.0 dataset (max/min ratio of 1.54 at L=1
and 2.00 at L=5, vs 1.14 and 1.39 for ont98). This is consistent with the
older basecaller having more difficulty with G/C homopolymers, a pattern that
Dorado v5.0 substantially ameliorates.

\subsection{Interface variant rate across basecaller versions}

The expanded interface analysis (Section~\ref{sec:interface}) repeated on
ont37 (651 specimens evaluated, with cluster recombination):

\begin{table}[h]
\centering
\caption{Interface vs non-HP variant rate ratio across both datasets.}
\label{tab:interface_all}
\begin{tabular}{lrlrl}
\toprule
Dataset & Specimens & Rate ratio & 95\% CI & Wilcoxon $p$ \\
\midrule
ont37 & 651 & 1.026 & 0.880--1.193 & 0.589 \\
ont98 & 663 & 0.880 & 0.785--0.988 & 0.042 \\
\bottomrule
\end{tabular}
\end{table}

Both datasets are consistent: interface HP positions do not show excess
variant rates relative to the non-HP baseline. The ont37 ratio of 1.026 has a
confidence interval spanning 1.0, indicating no detectable effect even with the
noisier v3.5 basecaller. The conclusion that interface positions require no
special CER treatment holds across basecaller generations.

\subsection{Base composition effects (ont98)}

We examined whether error rates need to be segmented by base (A/T/C/G) at each
HP length, using the ont98 filtered data:

\begin{table}[h]
\centering
\caption{Per-base variation in filtered HP error rates (ont98).}
\label{tab:base_effects}
\begin{tabular}{rlll}
\toprule
HP Length & Max base & Min base & Max/Min \\
\midrule
1 & 0.74\% (C) & 0.65\% (A) & 1.14 \\
2 & 0.94\% (T) & 0.72\% (C) & 1.31 \\
3 & 1.04\% (T) & 0.84\% (A) & 1.25 \\
4 & 1.19\% (G) & 0.94\% (A) & 1.26 \\
5 & 1.40\% (G) & 1.00\% (A) & 1.39 \\
6 & 2.30\% (G) & 0.89\% (A) & 2.57 \\
7 & 4.14\% (C) & 1.05\% (T) & 3.94 \\
\bottomrule
\end{tabular}
\end{table}

At L=1--5, the worst-case base is at most 1.4$\times$ the best---a modest
penalty for using a single rate. Importantly, the ``worst'' base is not
consistent (T at L=2--3, G at L=4--5), so there is no stable A/T vs C/G split
in this range. At L~$\geq 6$, C/G diverges sharply from A/T (the
well-established nanopore pattern), but sample sizes are small and these
lengths are approaching the limit of reliable variant calling regardless.


\section{Sensitivity to clustering identity threshold}
\label{sec:identity_sensitivity}

Speconsense's \texttt{-{}-min-identity} parameter controls which reads are
included in a cluster's MSA. At the default value of 0.90, reads with $<90\%$
pairwise identity to the cluster consensus are excluded. To test whether this
quality filter materially affects our error rate estimates, we repeated the
ont37 and ont98 analyses with \texttt{-{}-min-identity~0.85}.

\begin{table}[h]
\centering
\caption{Effect of \texttt{-{}-min-identity} on filtered per-position HP error
rates and non-HP error rates.}
\label{tab:identity_sensitivity}
\begin{tabular}{rrrrrcc}
\toprule
 & \multicolumn{2}{c}{ont37 (v3.5)} & \multicolumn{2}{c}{ont98 (v5.0)} & \multicolumn{2}{c}{085/default} \\
\cmidrule(lr){2-3} \cmidrule(lr){4-5} \cmidrule(lr){6-7}
HP Len & default & 0.85 & default & 0.85 & ont37 & ont98 \\
\midrule
1 & 1.56\% & 1.75\% & 0.67\% & 0.75\% & 1.12 & 1.11 \\
2 & 1.84\% & 2.05\% & 0.83\% & 0.92\% & 1.11 & 1.10 \\
3 & 2.17\% & 2.37\% & 0.97\% & 1.05\% & 1.10 & 1.08 \\
4 & 2.27\% & 2.52\% & 0.99\% & 1.06\% & 1.11 & 1.08 \\
5 & 2.13\% & 2.35\% & 1.13\% & 1.21\% & 1.10 & 1.07 \\
\midrule
\multicolumn{7}{l}{Non-HP total error rate} \\
 & 3.97\% & 4.67\% & 1.67\% & 1.88\% & 1.18 & 1.13 \\
\bottomrule
\end{tabular}
\end{table}

Relaxing the identity threshold inflates error rates by ${\sim}$10--12\% at
L=1--5, with a slightly larger effect on ont37 (where the older basecaller
produces more reads near the identity boundary). The inflation is uniform
across HP lengths, indicating that the additional reads are generally noisier
rather than selectively worse at homopolymers.

Crucially, the basecaller scaling factor is unaffected: the ont37/ont98 ratio
at \texttt{-{}-min-identity~0.85} is 2.33 at L=1 and 2.37 at L=4, compared to
2.32 and 2.30 at the default threshold. The identity filter and basecaller
version act as independent, multiplicative effects on the effective error rate.

For CER parameterization, the error rate should be estimated from data
processed with the same \texttt{-{}-min-identity} setting that will be used in
production. The results in this report use the default value of 0.90 unless
otherwise noted.


\section{Re-estimation under updated pipeline}
\label{sec:reestimation}

The original analysis (Sections~\ref{sec:interface}--\ref{sec:identity_sensitivity})
parameterized the shipped \texttt{dorado-v5.0.yaml} $q_\text{ctx}$ table used by
the context-aware CER framework in production~\cite{cer_practice}. The
speconsense pipeline has since gained several mechanisms---context-aware
variant phasing at $L \leq 5$ HP positions, cross-cluster read
reassignment, discard recovery, identity-group structure, and
\texttt{err\_factor}-based cluster homogeneity gating---each of which could
shift measured error rates. We re-estimated rates under the current pipeline
(speconsense 0.7.7, 2026-04-21) on the same ont98 input data to validate or
update the shipped values.

\subsection{Pipeline changes and expected directional effects}
\label{sec:pipeline_changes}

Table~\ref{tab:pipeline_changes} summarizes the relevant changes and their
expected effect on measured error rates. The changes fall into three classes:
cluster-cleaning (which should \emph{lower} apparent rates by producing
cleaner single-template primaries), noise-admitting (which should \emph{raise}
rates by bringing in reads that used to be rejected), and selection-redefining
(which change which specimens qualify for analysis without directly changing
rates).

\begin{table}[h]
\centering
\caption{Pipeline changes and their expected effect on measured HP/non-HP
error rates. ``$\downarrow$'' = expected to lower rates, ``$\uparrow$'' = raise,
``$\circ$'' = neutral.}
\label{tab:pipeline_changes}
\small
\begin{tabular}{lcl}
\toprule
Change & Direction & Mechanism \\
\midrule
Context-aware HP phasing ($L\leq5$) & $\downarrow$ & Splits HP-only paralogs into siblings \\
Cross-cluster read reassignment & $\downarrow$ & Reads move to best-matching cluster at variants \\
Tighter \texttt{-{}-min-variant-count}\ (5$\to$3) & $\downarrow$ & More aggressive non-HP phasing \\
Linear SPOA scoring (match~$+1$, else~$-1$) & $\downarrow$ & Unbiased alignment improves variant phasing \\
\texttt{err\_factor} cluster gating & $\downarrow$ & Selects self-consistent clusters \\
\midrule
Discard recovery & $\uparrow$ & Re-admits initially-discarded reads \\
\midrule
Identity group / complete linkage & $\circ$ & Affects CER decisions, not $q_\text{ctx}$ \\
Ambiguity calling (IUPAC) & $\circ$ & Few positions affected in practice \\
\bottomrule
\end{tabular}
\end{table}

\subsection{Methods}

ont98 (955 specimens) was re-processed with the current pipeline defaults,
preserving only the sampling configuration that matches the original
analysis:
\texttt{-{}-presample~0 -{}-max-sample-size~10000 -{}-collect-discards
-{}-error-model~dorado-v5.0}. All other parameters took current defaults
(\texttt{-{}-min-identity~0.9}, \texttt{-{}-group-identity~0.85},
\texttt{-{}-hp-normalization-length~6}, \texttt{-{}-min-variant-count~3},
\texttt{-{}-significance-level~1e-5}). The pipeline also uses a linear
SPOA scoring scheme (match~$+1$, mismatch~$-1$, gap~open~$-1$,
gap~extend~$-1$) in place of SPOA's defaults, which bias toward
substitutions over gaps. No failures.

\paragraph{Selection criterion evolution.} The original analysis selected
specimens whose primary cluster had $\geq 200$ reads and represented
$\geq 70\%$ of total reads (``primary\_frac'' criterion). Under the more
aggressive phasing of the current pipeline, this criterion becomes unworkable:
on ont98 it drops from 245 to 157 qualifying specimens because cluster
counts rise (median 10 clusters per specimen, with many small siblings
pulling the primary fraction below 0.70). We therefore switched to an
\texttt{err\_factor}-based criterion: the primary anchor cluster must have
\texttt{err\_factor}~$<~1.0$ and \texttt{RiC}~$\geq~200$. This yields 445
qualifying specimens on ont98. The criterion is peer-independent (does not
depend on sibling structure), directly measures single-template plausibility,
and is stable under the current pipeline.

In practice the \texttt{err\_factor} gate is non-binding at
\texttt{RiC}~$\geq~200$: all 445 qualifying ont98 specimens already
satisfy \texttt{err\_factor}~$<~1.0$ (maximum observed value 0.929).
\texttt{RiC} is therefore the effective selection criterion;
\texttt{err\_factor} serves here as a backup guarantee rather than an
active filter. This distinction matters for the circularity argument:
because $q_\text{ctx}$ enters the selection only through
\texttt{err\_factor}, and \texttt{err\_factor} is not constraining our pool,
the rate estimates are not circularly dependent on the $q_\text{ctx}$ value
used to compute them. Section~\ref{sec:revised_qctx} confirms this
analytically under the proposed table revision.

\paragraph{Approach 2: all-cluster non-HP analysis.} To characterize the
operational non-HP rate CER evaluates against in production (which includes
candidate clusters of all sizes), we additionally ran a weighted analysis
across all clusters with \texttt{RiC}~$\geq~5$, using each cluster's own
consensus as its local scaffold. Results are stratified by \texttt{RiC} bin.

\subsection{Results: rate reduction under the updated pipeline}

Table~\ref{tab:reestimation_summary} compares the original filtered rates
(primary\_frac~$\geq~0.70$ criterion on old-pipeline clusters) to the re-
estimated rates (\texttt{err\_factor}~$<~1.0$ criterion on new-pipeline
clusters). Across all contexts, measured rates are consistently lower under
the current pipeline, by a factor of approximately 0.70--0.80.

\begin{table}[h]
\centering
\caption{Filtered per-position error rates on ont98: original analysis vs.\
re-estimation under the current pipeline. ``v2\_ef'' denotes the current-
pipeline rate under the \texttt{err\_factor}~$<~1.0$ selection.}
\label{tab:reestimation_summary}
\begin{tabular}{lrr}
\toprule
Context & original & v2\_ef \\
\midrule
non-HP sub          & 0.0059 & 0.0039 \\
non-HP del          & 0.0040 & 0.0028 \\
non-HP ins          & 0.0068 & 0.0046 \\
\midrule
hp-l1               & 0.0067 & 0.0047 \\
hp-l2               & 0.0083 & 0.0062 \\
hp-l3               & 0.0097 & 0.0076 \\
hp-l4               & 0.0099 & 0.0081 \\
hp-l5               & 0.0113 & 0.0087 \\
\bottomrule
\end{tabular}
\end{table}

\paragraph{Framing observation.} The current pipeline admits \emph{more}
reads into clusters than the original (via discard recovery and cross-
cluster read reassignment). The a priori expectation under such additions
is that per-cluster error rates should \emph{rise}: more reads means
more chances to admit noisier reads into any given cluster. The observed
direction is the opposite---rates fall by 0.70--0.80$\times$. The most
parsimonious explanation is that the new cluster-composition mechanisms
(reassignment and recovery placing reads into their best-matching clusters
based on variant positions) produce cleaner per-cluster read populations
than the original simpler MCL-only assignment, and that this cleaning
effect dominates any noise admitted by expanding the read pool.

\subsection{All-cluster non-HP rates by RiC bin}

Table~\ref{tab:approach2} reports the all-cluster pooled non-HP rates
stratified by \texttt{RiC}. On ont98, rates are essentially flat across the
\texttt{RiC} 20--1000+ range (substitution rate varies from 0.0036 to 0.0039).
The \texttt{RiC} 5--20 bin is slightly elevated ($\sim$20\%). The pooled
non-HP rate (0.0039) is within 3\% of the primary-anchor-only rate (0.0038),
indicating that admitting small clusters does not materially shift the
aggregate estimate.

\begin{table}[h]
\centering
\caption{All-cluster non-HP error rates by \texttt{RiC} bin (ont98). Each row
pools reads across all MSAs (\texttt{RiC}~$\geq~5$) in the specimen, using
per-cluster consensus as the ground-truth scaffold.}
\label{tab:approach2}
\small
\begin{tabular}{lrrr}
\toprule
Bin (\texttt{RiC}) & Clusters & Sub & Sub+Del \\
\midrule
5--20             & 5,472 & 0.0046 & 0.0078 \\
20--50              & 632 & 0.0036 & 0.0058 \\
50--100             & 324 & 0.0039 & 0.0063 \\
100--200            & 289 & 0.0037 & 0.0062 \\
200--500            & 408 & 0.0039 & 0.0066 \\
500--1000           & 203 & 0.0038 & 0.0066 \\
1000--$\infty$        & 29 & 0.0037 & 0.0065 \\
\midrule
POOLED           & 7,357 & 0.0039 & 0.0066 \\
\bottomrule
\end{tabular}
\end{table}

\subsection{Mechanism: isolating HP phasing contribution (Test A)}
\label{sec:test_a}

A natural concern is that the more aggressive HP phasing at $L \leq 5$ in
the current pipeline could produce \emph{artifactual} splits---cutting noise
from signal and leaving a cleaner-than-reality primary cluster, thereby
biasing apparent rates downward. To test this, we re-ran ont98 with
\texttt{-{}-hp-normalization-length~0}, which reverts to the old behavior of blanket-
normalizing all HP variants (no HP-driven splits).

Table~\ref{tab:test_a} compares HP rates between the current-pipeline run
(\texttt{v2\_ef}) and the blanket-normalization variant (Test~A). Rates
shift by at most 0.0003 at $L \leq 5$. If HP phasing were producing
substantial artifactual splits, disabling it would recover a large fraction
of the gap to the original-pipeline rates; the observed gap recovery is
only 2--11\%.

\begin{table}[h]
\centering
\caption{Effect of disabling HP phasing (\texttt{-{}-hp-normalization-length~0}) on
measured HP rates. Gap recovery is the fraction of the v2\_ef$\to$original
gap closed by Test~A.}
\label{tab:test_a}
\begin{tabular}{lrrrrr}
\toprule
L & v2\_ef & Test~A raw & Test~A filt & filt delta & gap recovery \\
\midrule
1 & 0.0047 & 0.0049 & 0.0048 & +0.0000 & 2\% \\
2 & 0.0062 & 0.0064 & 0.0062 & +0.0000 & 2\% \\
3 & 0.0076 & 0.0080 & 0.0076 & +0.0001 & 3\% \\
4 & 0.0081 & 0.0087 & 0.0082 & +0.0001 & 6\% \\
5 & 0.0087 & 0.0096 & 0.0090 & +0.0003 & 11\% \\
\bottomrule
\end{tabular}
\end{table}

Two side-channel measurements confirm that HP phasing \emph{is} doing real
work, but that its output is absorbed by downstream filters:

\begin{itemize}
\item Cluster counts shifted under Test~A: 387/955 specimens produced fewer
clusters (HP-only splits now merged), 521 unchanged, 47 produced more. HP
phasing is not inert; it is doing real partitioning.
\item Bimodal flag counts (Step~6b of the biological-variant filter)
doubled under Test~A (22$\to$46). When HP variants are merged back into the
primary, they surface as bimodal length distributions, which the filter
excludes from the filtered rate. The outlier-detection filter also caught an
additional 122 positions.
\end{itemize}

The \emph{filtered} rate aggregation is therefore robust to HP phasing
behavior: merging HP variants back in doubles the bimodal flag count, the
filter catches them, and the resulting per-context rates are nearly
unchanged. This establishes that the bimodal and outlier filters
(Section~\ref{sec:interface}, \emph{Biological variant filtering}) are a
structural safety net for rate measurement, not a cosmetic post-hoc
adjustment.

\paragraph{Implication for the rate drop.} HP phasing accounts for
approximately 2--11\% of the measured rate reduction between the original
and current pipelines. The remainder (${\sim}90\%$) comes from other
changes, principally:

\begin{itemize}
\item \textbf{Cross-cluster read reassignment} (likely dominant): reads that
were initially placed in a suboptimal cluster are moved to the cluster whose
consensus they actually match at variant positions. Apparent per-read
disagreement within each cluster falls accordingly. This mechanism changes
cluster \emph{membership}, not structure, so it is not isolated by Test~A.
\item \textbf{Linear SPOA scoring} (match~$+1$, mismatch~$-1$,
gap~open~$-1$, gap~extend~$-1$) in place of SPOA's defaults. The
defaults bias toward substitutions over gaps, which forces reads with
real indels into substitution-heavy alignments against the consensus
and blurs the signal that variant phasing uses to separate paralogs.
The linear scheme treats substitutions and indels equivalently, so
indel-distinguished haplotypes phase cleanly and the resulting
primary clusters are more homogeneous.
\item More aggressive non-HP phasing (\texttt{-{}-min-variant-count}\ 5$\to$3)
separating non-HP paralogs into siblings.
\item Selection criterion change (\texttt{primary\_frac} to
\texttt{err\_factor}) admitting a different specimen population.
\end{itemize}

Distinguishing among these contributions would require further test runs
with individual flags reverted. For the purpose of this re-estimation, the
relevant conclusion is that the observed rates are \emph{not} inflated by
artifactual HP splitting.

\subsection{Revised shipped $q_\text{ctx}$ table}
\label{sec:revised_qctx}

Table~\ref{tab:revised_qctx} gives the proposed update to
\texttt{dorado-v5.0.yaml}. HP rates use the \texttt{v2\_ef} primary-anchor
values (approach-1 regime, where mode-based ground truth is most reliable).
Non-HP rates use the all-cluster pooled values from approach 2, which match
the operational distribution CER evaluates against.

\begin{table}[h]
\centering
\caption{Proposed revised $q_\text{ctx}$ values for
\texttt{dorado-v5.0.yaml}. ``Original'' values are the current shipped
table (from the original analysis); ``Revised'' are the re-estimated values
from this section.}
\label{tab:revised_qctx}
\begin{tabular}{lrrc}
\toprule
Context & Original & Revised & Basis \\
\midrule
\texttt{non-hp-sub}   & 0.0059 & 0.0039 & approach 2 pooled (ont98) \\
\texttt{non-hp-indel} & 0.0108 & 0.0074 & approach 2 pooled (ont98) \\
\midrule
\texttt{hp-l1}        & 0.0067 & 0.0047 & v2\_ef primary (ont98) \\
\texttt{hp-l2}        & 0.0083 & 0.0062 & " \\
\texttt{hp-l3}        & 0.0097 & 0.0076 & " \\
\texttt{hp-l4}        & 0.0099 & 0.0081 & " \\
\texttt{hp-l5}        & 0.0113 & 0.0087 & " \\
\bottomrule
\end{tabular}
\end{table}

The revised values shift \texttt{cer\_factor} outputs upward by
approximately $1/0.72 \approx 1.39\times$ on existing data. No upstream
re-runs are required to adopt the new table; per the metadata-2.0 schema
design~\cite{cer_practice}~\S6, stored output can be recomputed from the
retained $N$, $M$, $K$, and context tags.

We confirmed analytically that the revised table does not induce a selection
shift at our analysis tier: \texttt{err\_factor} values scale approximately
as $1/q_\text{ctx}$, so the revision inflates primaries'
\texttt{err\_factor} by ${\sim}1.39\times$. The maximum observed primary
\texttt{err\_factor} in the current re-estimation (0.929) rises to
$\approx 1.29$ under the revised table---still below the default
\texttt{-{}-max-err-factor~1.5} filter threshold. Selection is therefore
stable; no bootstrap iteration is required.

\subsection{ont37 retune (Dorado v3.5)}
\label{sec:reestimation_ont37}

The same procedure was applied to ont37 to refresh
\texttt{dorado-v3.5.yaml}. Re-clustering was performed under speconsense
0.8.1 (released after the v5.0 retune) with
\texttt{-{}-error-model~dorado-v3.5} (chemistry-matched to avoid
underestimating the noise floor).

\paragraph{Selection criterion.} ont37's lower per-specimen depth (median
primary RiC was 59 in the original analysis) interacts with the more
aggressive context-aware phasing of 0.8.x: primaries shrink further, and
\texttt{RiC}~$\geq~200$ yields only 10 specimens. We relaxed the floor to
\texttt{RiC}~$\geq~50$, yielding 80 qualifying specimens (vs.\ 94 under the
original \texttt{RiC}~$\geq~100$, \texttt{primary\_frac}~$\geq~0.50$
criterion). The \texttt{err\_factor}~$<~1.0$ gate is empirically
\emph{fully} non-binding at this threshold (maximum observed value 0.737),
preserving the circularity defense for the q\_ctx-driven selection. The
revised table inflates this maximum to $\approx 1.32$ under
$1/q_\text{ctx}$ scaling---still below the \texttt{-{}-max-err-factor~1.5}
cap---so selection remains stable after the table refresh.

\paragraph{Results.} Table~\ref{tab:reestimation_ont37} compares the
original ont37 rates to the v2\_ef re-estimation. Rate ratios are broadly
similar to ont98's (0.56--0.74$\times$ on ont37 vs.\ 0.66--0.82$\times$ on
ont98), confirming that the pipeline mechanisms responsible for the v5.0
rate drop apply to v3.5 data as well. The v3.5/v5.0 ratio narrows somewhat
at higher $L$ under the new pipeline (2.23$\times$ at $L=1$ down to
1.41$\times$ at $L=5$, vs.\ a more uniform $\sim 2.2\times$ under the
original) but remains well above unity at every short HP length.

\begin{table}[h]
\centering
\caption{Filtered per-position error rates on ont37: original analysis vs.\
re-estimation under speconsense 0.8.1 (80 qualifying specimens at
\texttt{RiC}~$\geq~50$, \texttt{err\_factor}~$<~1.0$). Non-HP rates from
approach 2 (all-cluster pooled, \texttt{RiC}~$\geq~20$, 184 clusters).}
\label{tab:reestimation_ont37}
\begin{tabular}{lrrr}
\toprule
Context & Original & v2\_ef & Ratio \\
\midrule
\texttt{non-hp-sub}   & 0.0144 & 0.0081 & 0.56 \\
\texttt{non-hp-indel} & 0.0222 & 0.0158 & 0.71 \\
\midrule
\texttt{hp-l1}        & 0.0156 & 0.0105 & 0.67 \\
\texttt{hp-l2}        & 0.0184 & 0.0137 & 0.74 \\
\texttt{hp-l3}        & 0.0217 & 0.0150 & 0.69 \\
\texttt{hp-l4}        & 0.0227 & 0.0142 & 0.63 \\
\texttt{hp-l5}        & 0.0213 & 0.0123 & 0.58 \\
\bottomrule
\end{tabular}
\end{table}

\paragraph{Caveats specific to ont37.} The qualifying pool is smaller than
ont98's (80 vs.\ 445 specimens) and has lower per-specimen depth
(\texttt{RiC} range 51--396 vs.\ ont98's $\geq 200$). HP run counts at
$L=4$--$5$ are correspondingly limited (526 and 154 runs respectively),
producing slight non-monotonicity in the rate-vs-length curve at high $L$.
The bimodal/outlier filter exclusion rate is low (0.5\%, 217/43,276 runs),
indicating the filter's biological-variant cleanup is operating
conservatively at this specimen count. These values should therefore be
treated as the best current estimates rather than as tightly constrained
measurements; the test in
\texttt{tests/test\_qctx.py:test\_v3\_5\_rates\_higher\_than\_v5\_0\_at\_short\_lengths}
remains satisfied (v3.5 rates exceed v5.0 by 1.75$\times$--2.23$\times$ at
$L=1$--$4$).

\paragraph{Convergence (fixed-point check).} Because the shipped table
participates in the clustering step via the CER significance gate
(through \texttt{err\_factor} and variant phasing decisions), we verified
that ont37 has converged to a fixed point of the bootstrap. ont37 was
re-clustered a second time under the revised \texttt{dorado-v3.5.yaml}
and the analysis re-run. Selection was highly stable (80 specimens both
passes, 77 shared), and measured rates drifted by at most 2.5\%
relative: \texttt{hp-l1}~$-2.40\%$, \texttt{hp-l2}~$-2.49\%$,
\texttt{hp-l3}~$-1.46\%$, \texttt{hp-l4}~$+0.41\%$, \texttt{hp-l5}~$+0.36\%$;
\texttt{non-hp-sub}~$-0.61\%$, \texttt{non-hp-indel}~$-0.66\%$. The mild
downward direction at $L=1$--$3$ is mechanistic: under the lower q\_ctx
expectations, the CER gate is more permissive at HP positions, so a few
additional variants phase out into siblings and primaries are
infinitesimally cleaner. All drifts sit within the sampling noise of the
80-specimen pool, so the shipped values are retained without further
iteration. Median primary \texttt{err\_factor} inflated from 0.541 to
0.731 between passes (less than the $1/q_\text{ctx}$-only prediction of
$\sim 1.78\times$, because cluster cleanup under the new yaml partially
offsets the denominator drop); maximum observed value 1.092 remains well
below the \texttt{-{}-max-err-factor~1.5} cap.

\subsection{Limitations}

\begin{itemize}
\item \textbf{ITS-specific.} The re-estimation inherits the amplicon-specific
limitation of the original analysis. Other amplicons or organism contexts
may produce different values.
\item \textbf{iNat input distribution.} The measurements reflect specimen
quality distributions typical of iNaturalist uploads (ont98). Environmental
DNA or specimen sources with substantially different input-quality variance
could produce different all-cluster rates.
\item \textbf{Run-to-run variation is not characterized in this report.}
The shipped table is derived from a single ont98 run; runs of different
provenance (other Mycota Lab batches, other operators) may produce
systematically different rates. The CER factor's linear sensitivity to
$q_\text{ctx}$~\cite{cer_practice}~\S3 handles this by design.
\item \textbf{Mechanism attribution is not fully resolved.} Test~A
eliminates HP phasing as the dominant contributor. Definitive identification
of the dominant mechanism (cross-cluster reassignment vs.\ non-HP phasing
refinements vs.\ selection criterion) would require additional test runs
with individual changes reverted; these are future work.
\item \textbf{Phase-2 per-run estimation} (\cite{cer_practice}~\S4.2) would
eventually replace the shipped table entirely, eliminating the
basecaller/chemistry-update dependency. The revised values are an interim
update pending that work.
\end{itemize}

\section{Discussion}

\subsection{Recommended approach for speconsense}

The combined results from two sequencing runs spanning two basecaller
generations support a simple two-tier approach:

\textbf{Short homopolymers (length 1--5):} Implied per-position error rates
are consistently below 2.3\% (filtered) across both datasets, all bases, and
both basecaller versions. For the current Dorado v5.0 basecaller, rates are
below 1.4\%. These positions can be brought into the CER framework---no
special segmentation by base is needed (max/min ratio across bases is only
1.1--1.4$\times$ at these lengths for v5.0, and 1.2--2.0$\times$ for v3.5).
The CER implementation should evaluate hp variants at the run level (i.e.,
``what is the probability of observing hp length $k$ when the true length is
$L$?'') using the empirical error distributions from this analysis, rather than
treating each position in the run as an independent trial. The per-position
rate serves as a compact parameterization of the run-level error probability
for a given (base, length) context, but the significance test itself operates
on the run as a unit. The companion report~\cite{cer_practice} develops this
into a deployed framework: per-basecaller tables of $q_\text{ctx}(\text{base},
\text{length})$ derived from this analysis feed the binomial survival test, and
significance is reported as a dimensionless \texttt{cer\_factor} (the
ratio of the critical joint error rate to the product of empirical $q_\text{ctx}$
values) that remains comparable across contexts and re-scales linearly as tables
are refreshed for new chemistry or basecaller generations. The initial production
deployment uses shipped per-basecaller tables pooled across base and direction;
per-run estimation to account for flow-cell variation
(\cite{cer_practice}~\S 4.2) is left as future work.

\textbf{Long homopolymers (length 6+):} Error rates rise steeply and diverge
significantly by base (C/G 2--4$\times$ worse than A/T at L=6--7). Sample
sizes are small, and the error rates are high enough that variant calling
becomes unreliable. The current blanket normalization approach (treating all HP
length differences as insignificant) remains appropriate for these lengths.
Enhanced positional error reporting in the quality report could flag positions
where the observed HP length distribution suggests a potential variant, even if
the CER framework cannot confidently evaluate it.

\textbf{Interface positions:} No special treatment required. Well-supported
variants at HP interface positions occur at or below the rate of non-HP
variants, confirmed across both datasets (rate ratio 1.03 in ont37, 0.88
in ont98).

\subsection{Interface homopolymers}

The finding that interface hp substitutions are nearly as prevalent as pure hp
length errors (0.97\% vs 1.06\%) initially raised the question of whether
interface substitutions might be signal-level artifacts (boundary confusion in
the nanopore signal). The expanded analysis resolves this question decisively
across both datasets:

\begin{enumerate}
\item Well-supported variants at HP interface positions occur at or below the
  rate of non-HP variants (ont37 ratio 1.03 [95\% CI: 0.88--1.19]; ont98 ratio
  0.88 [95\% CI: 0.79--0.99]). There is no excess that would indicate a
  signal-level artifact, even with the noisier v3.5 basecaller.

\item Among interface variant positions, the fraction with an
  interface-compatible top allele is close to the 33\% expected by chance.
  Boundary-confusion substitutions are not enriched.

\item No published study has specifically analyzed substitution error rates at
  HP boundaries in nanopore data. Apparent boundary effects reported in the
  literature (e.g., Liu-Wei et al.~\cite{liuwei2024}, Dunn et
  al.~\cite{dunn2023}) are alignment artifacts from minimap2 representing HP
  length errors as SNPs at run boundaries, not signal-level errors.
\end{enumerate}

\subsection{Comparison to literature}

Both datasets show substantially better HP accuracy than published R10.4
/ Guppy SUP values~\cite{sereika2022}, with progressive improvement across
basecaller versions. Note that this comparison is approximate: Sereika measured
raw read-level accuracy against a known reference genome, whereas our rates are
effective error rates conditioned on speconsense's post-clustering MSA
(Section~2.5). The quality filtering inherent in clustering excludes the
noisiest reads, so our effective rates are expected to be somewhat lower than
raw read-level rates for the same chemistry and basecaller. The magnitude of
the difference attributable to this methodological difference vs.\ genuine
basecaller improvement cannot be separated without matched raw-read
measurements on our data.

\begin{table}[h]
\centering
\caption{Whole-run accuracy progression from Sereika (Guppy) through our
datasets (Dorado v3.5 and v5.0). Our values are effective post-clustering
rates (Section~2.5), not directly comparable to Sereika's raw read-level
measurements.}
\begin{tabular}{rrrr}
\toprule
HP Len & Sereika R10.4 & ont37 (v3.5) & ont98 (v5.0) \\
\midrule
1 & ${\sim}$99.2\% & 98.4\% & 99.3\% \\
3 & ${\sim}$94.5\% & 93.6\% & 97.1\% \\
5 & ${\sim}$86.0\% & 89.8\% & 94.5\% \\
7 & ${\sim}$60.0\% & 74.0\% & 85.4\% \\
\bottomrule
\end{tabular}
\end{table}

The ont37 data (Dorado v3.5) provides an intermediate point between the
published Guppy results and our v5.0 data, demonstrating a clear trajectory
of improvement. Even at the v3.5 level, short HP accuracy is sufficient for
CER analysis. The cross-basecaller comparison reveals a substantial
version-level effect (${\sim}$2.2$\times$ from v3.5 to v5.0). Run-to-run
variation within a basecaller version is plausible (flow cell quality,
library preparation, basecaller calibration drift) but is not characterized
in this report; per-run error rate estimation in production is left as
future work.

\subsection{Limitations and open questions}

\begin{enumerate}
\item \textbf{Pipeline-conditioned measurement}: Our error rates are effective
  rates at the post-clustering stage, conditioned on speconsense's identity
  threshold and clustering parameters (Section~2.5). They are not directly
  comparable to raw read-level error rates in the literature. The
  ${\sim}$10--12\% sensitivity to the identity threshold
  (Section~\ref{sec:identity_sensitivity}) is modest but means that CER
  parameters should be estimated from data processed with the same settings
  used in production.

\item \textbf{Run-specific error rates}: The ${\sim}$2.2$\times$ scaling
  factor between basecaller versions (v3.5 to v5.0) suggests that HP error
  rates should be estimated per-basecaller rather than assumed universal.
  Within-basecaller run-to-run variation (flow cell quality, library
  preparation, basecaller calibration drift) is expected but is not
  characterized in this report.

\item \textbf{Ground truth}: We use the empirical mode as ground truth, which
  is reliable for hp $\leq {\sim}9$ at our read depths but may be incorrect for
  longer runs. Cross-platform validation with PacBio HiFi data would address
  this.

\item \textbf{PCR artifacts}: PCR can introduce hp length errors that are
  physically present in the library. These are indistinguishable from sequencing
  errors in our analysis and are out of scope for the CER framework (which
  establishes physical library presence, not biological reality).

\item \textbf{Sample composition}: Our data is exclusively ITS amplicons from
  fungi. The hp length distribution and context effects may differ for other
  loci or organisms, though the physical error mechanism should be universal.

\item \textbf{Biological variant filtering}: The outlier/bimodal detection
  thresholds are heuristic. A more sophisticated approach (e.g., cross-specimen
  consistency checks) could improve filtering, though the current approach is
  sufficient for the cross-basecaller comparison.
\end{enumerate}


\section*{Acknowledgments}

All sequencing data analyzed in this report were produced by Mycota Lab
(Hoosier Mushroom Society) under the protocol of
Russell~\cite{russell2023}. The author thanks Stephen D.~Russell for sharing
the run-level outputs that made this analysis possible.

Substantial portions of the text were drafted with the assistance of Claude
(Anthropic, Opus 4.6). The core ideas, analytical framing, and mathematical
approach were provided by the author; the AI contributed to exposition,
\LaTeX{} formatting, and iterative drafting under human direction. The author
takes full responsibility for the content.


\begin{thebibliography}{99}

\bibitem{delahaye2021}
C.~Delahaye and J.~Nicolas.
\newblock Sequencing {DNA} with nanopores: Troubles and biases.
\newblock \emph{PLoS ONE}, 16(10):e0257521, 2021.

\bibitem{dunn2023}
T.~Dunn, E.~Berry, C.~M. Bergman.
\newblock {nPoRe}: $n$-polymer read correction for improved nanopore
homopolymer accuracy.
\newblock \emph{BMC Bioinformatics}, 24:292, 2023.

\bibitem{liuwei2024}
W.~Liu-Wei, W.~Trzaskowski, et~al.
\newblock Benchmarking {Oxford Nanopore} long-read {RNA} sequencing for
transcript quantification.
\newblock \emph{BMC Genomics}, 25:219, 2024.

\bibitem{russell2023}
S.~D. Russell.
\newblock {ONT DNA} barcoding fungal amplicons w/ {MinION \& Flongle V.3}.
\newblock \emph{protocols.io}, March 2023.
\newblock \url{https://dx.doi.org/10.17504/protocols.io.36wgq7qykvk5/v3}.

\bibitem{ni2023}
Y.~Ni, G.~Liu, H.~Peng, et~al.
\newblock Benchmarking of {Nanopore} {R10.4} and {R9.4.1} flow cells in
single-cell whole-genome amplification and whole-genome shotgun sequencing.
\newblock \emph{Computational and Structural Biotechnology Journal}, 21:2352--2363, 2023.

\bibitem{sereika2022}
M.~Sereika, R.~H. Kirkegaard, S.~M. Karst, et~al.
\newblock Oxford {Nanopore} {R10.4} long-read sequencing enables the generation
of near-finished bacterial genomes from pure cultures and metagenomes without
short-read or reference polishing.
\newblock \emph{Nature Methods}, 19:482--488, 2022.

\bibitem{walker2026}
J.~Walker.
\newblock Critical error rate analysis for {MSA}-based variant phasing in
nanopore amplicon sequencing.
\newblock Technical report, 2026.
\newblock \url{https://github.com/joshuaowalker/speconsense/blob/main/docs/variant_significance/variant_significance.pdf}

\bibitem{cer_practice}
J.~Walker.
\newblock Critical error rate in practice: context-aware variant significance
for nanopore amplicon pipelines.
\newblock Technical report, 2026.
\newblock \url{https://github.com/joshuaowalker/speconsense/blob/main/docs/cer_in_practice/cer_in_practice.pdf}

\bibitem{wick2025}
R.~R. Wick.
\newblock Benchmarking {Oxford Nanopore} bacterial genome assembly.
\newblock \emph{Microbial Genomics}, 2025.

\end{thebibliography}

\end{document}
